% =====================================================================
%  CHAPTER 5: 儿童心理发育（对应大纲第五章）
%  内容来源：往届笔记第四章，参考PPT第四章
% =====================================================================
\chapter{儿童心理发育}
\setcounter{chapter}{5}
\chapterintro{
\textbf{掌握：}儿童心理发育的特点：儿童少年认知能力发育、儿童少年情绪情感发育、儿童少年个性及社会化发展。
}

\begin{defbox}
\textbf{心理行为发育（psychological and behavior development）：}从受精卵开始到出生、成熟，儿童心理和行为特征的变化过程，包括认知、语言、情绪情感、人格和社会化适应等多个方面。
\end{defbox}

% ===== 掌握内容 =====
\section{儿童少年认知能力发育}

\begin{defbox}
\textbf{认知（cognition）：}认识活动或认识过程，是大脑反映客观事物的特征、状态及其相互关系、并揭示事物对人的意义与作用的一种高级心理活动。
\end{defbox}

\subsection{感知觉发育}

\begin{defbox}
\textbf{感知觉（sensory perception）：}人脑对当前作用于感觉器官的客观事物的反映。

\textbf{感觉（sense）：}对事物个别属性的反映，依赖于个别感觉器官的活动，分为视觉、听觉、嗅觉、味觉和触觉。

\textbf{知觉（perception）：}人对事物整体属性的综合反映，是在感觉基础上发展起来的，依赖多种感觉器官，分为时间知觉、空间知觉和运动知觉。
\end{defbox}

\begin{keybox}
\textbf{一、视觉发育}

\begin{enumerate}[leftmargin=*]
    \item 2个月：视觉集中明显形成（尤其是人脸）。
    \item 4个月：出现辨色视觉。
    \item 6个月以前：视力发展敏感期，视力和立体视觉都逐渐发育。
    \item 2-3岁：能正确辨别红、黄、绿、蓝4种基本颜色并出现双眼视觉。
    \item 出生后3年内：视觉系统功能基本发育成熟。
\end{enumerate}

\textbf{二、听觉发育}

\begin{enumerate}[leftmargin=*]
    \item 出生时：听觉器官发育基本成熟，新生儿能听见声音及区分声音的频率、强度和持续时间。
    \item 2岁以后：听力已经很灵敏，几乎达到成人水平。
    \item 学龄初期儿童：听觉感受性随着年龄增长不断发展，在音调辨别、言语听觉、语音听觉和听觉敏感等方面明显提高。
    \item 青少年时期：听觉能力已基本达到稳定水平。
\end{enumerate}

\textbf{三、嗅觉发育}

\begin{enumerate}[leftmargin=*]
    \item 新生儿：嗅觉与成人一样敏感，出生时嗅觉发育已比较成熟，能表现出对不同气味的反应。
    \item 7-8个月：嗅觉逐渐灵敏，能分辨出芳香气味。
    \item 2岁左右：已能很好辨别各种气味。
\end{enumerate}

\textbf{四、味觉发育}

\begin{enumerate}[leftmargin=*]
    \item 味觉在儿童时期最发达，以后逐渐衰退。
    \item 2h的新生儿：能分辨甜、酸、苦、咸等多种味道。
    \item 6个月-1岁：幼儿在这一阶段味觉发展最灵敏。
\end{enumerate}

\textbf{五、触觉发育}

\begin{enumerate}[leftmargin=*]
    \item 触觉是人体发展最早、最基本的感觉；触觉在胎儿期已开始发育，到新生儿阶段已高度敏感，触觉是婴儿认识世界的主要手段之一。
    \item 对婴儿的抚触就是通过对触觉的刺激，增强婴儿触觉的敏感性，加强对外界的反应，促进其发育。
    \item 2-3岁：能很好辨别各种物体的不同属性。
    \item 5-6岁：皮肤觉分辨的精细度逐渐提高，在学龄前期已逐步趋于完善。
\end{enumerate}

\textbf{六、知觉发育}

\begin{enumerate}[leftmargin=*]
    \item 婴儿3个月时有了分辨简单形状的能力，6个月以前能辨别大小，9-12个月能知觉形状。
    \item 随着感觉功能的完善促使幼儿的感知觉进一步发展，主要表现为在空间知觉上辨别形状的能力逐渐增强。
    \item 学龄期儿童的视觉感受性和听觉感受性会随年龄增长而不断发展，空间知觉、方位知觉、距离知觉和时间知觉都有很大进步。
\end{enumerate}
\end{keybox}

\subsection{注意力和记忆力发育}

\begin{keybox}
\textbf{一、注意（attention）}

心理活动对一定对象有选择的集中，是认知过程的开始，分为无意注意和有意注意。

\begin{enumerate}[leftmargin=*]
    \item \textbf{学龄期儿童：}有意注意进一步发展，注意的稳定性随年龄增长而增强，注意集中时间也随之提高。
    \item \textbf{青少年期：}注意的集中性和稳定性不断增长，稳定性大概能维持40min。
\end{enumerate}

\textbf{二、记忆（memory）}

人脑对过去经验的反映，包括识记、保持和再现（回忆和再认），需经历感知觉记忆、短时记忆和长时记忆3个阶段。

\begin{enumerate}[leftmargin=*]
    \item \textbf{婴幼儿时期：}记忆迅速发展的第一个时期，条件反射的出现是记忆发生的标志。记忆容量增加显著，但记忆以无意识记为主，有意识记成分在逐渐增加。
    \item \textbf{学龄期：}由于系统学习需要记住大量东西，此时记忆发生质的变化。
    \item \textbf{青少年期：}记忆的整体水平处于人生的最佳时期。有意识记日益占主导地位，对抽象材料的记忆能力也明显增强。
\end{enumerate}
\end{keybox}

\subsection{思维和想象能力发育}

\begin{keybox}
\textbf{一、思维（thinking）}

人脑对客观事物间接、概括的反映，能认识事物的本质和事物间的内在联系，属认知的高级阶段。

\begin{enumerate}[leftmargin=*]
    \item \textbf{婴幼儿期：}思维依靠感知觉和动作完成；思维主要是直觉行动思维，是指婴幼儿在动作中进行的思维。婴幼儿在进行这种思维时，只能反映自身动作所触及的具体事物，而不能离开动作，在感知和动作外思考。
    \item \textbf{学龄前期：}直观行动思维→具体形象思维→抽象逻辑思维，其中占主导地位的是具体形象思维。
    \item \textbf{学龄期：}思维的基本特点是从具体形象思维逐步过渡到抽象逻辑思维，这是思维发展过程中质的变化。
    \item \textbf{青春期：}思维能力有很大发展，抽象逻辑思维处于优势地位，能运用抽象思维突破心理运算的界限和理解各种抽象概念，并获得更多增长新知识的机会。
\end{enumerate}

\textbf{二、想象（imagine）}

对已有表象进行加工改造，形成新形象的过程。思维是想象的基础。

\begin{enumerate}[leftmargin=*]
    \item \textbf{婴幼儿期：}1-2岁的婴幼儿出现想象的萌芽，主要通过动作和口头言语表达出来；3岁时，随着经验与言语的发育，逐渐产生了具有简单想象的游戏。
    \item \textbf{学龄前期：}已具有丰富的想象力，集中表现在游戏活动中，非游戏时想象水平较低。
    \item \textbf{学龄期：}想象的有意性、目的性逐渐增强，创造性成分逐渐增大，现实性逐渐提高。
    \item \textbf{青少年期：}想象能力随教学活动的深入进一步发展，主要表现为有意想象占主要地位，创造想象日益占优势地位，再创想象能力更加独立、概括和精确，想象内容更加复杂，抽象概括性和逻辑性更高。
\end{enumerate}
\end{keybox}

\subsection{语言发育}

\begin{defbox}
\textbf{语言（language）：}人类社会中客观存在的现象，是一种社会上约定俗成的符号系统，是由词汇（包括音形义）按一定语法构成的。
\end{defbox}

\begin{tipbox}
\textbf{语言发育阶段：}

\begin{enumerate}[leftmargin=*]
    \item \textbf{前语言阶段（prespeech stage）：}从婴儿出生到第一个有真正意义的词产生前的这一时期。
    \item \textbf{语音的发展：}学龄前儿童发音的正确率随年龄的增长而提高，对韵母的发音比较容易掌握，发音错误多出现于辅音中的翘舌音和齿音。
\end{enumerate}
\end{tipbox}

\section{儿童少年情绪情感发育}

\begin{defbox}
\textbf{情绪（emotion）：}客观事物是否符合人的需要产生的态度体验，反映客观事物与人的需要间的关系。

\textbf{情感（emotional feeling）：}与人的高级社会性需要满足与否相联系的态度体验，是在社会交往的实践中逐渐形成的，如友谊感、道德感、美感和理智感等，是人类独有的一种情绪状态。
\end{defbox}

\begin{keybox}
\textbf{儿童青少年情绪情感发育特点：}

\textbf{1. 婴幼儿情绪情感发育特点：}
\begin{enumerate}[leftmargin=*]
    \item 与生理需求是否得到满足直接相关。
    \item 是儿童与生俱来的遗传本能，有先天性。
\end{enumerate}

\textbf{2. 学龄前儿童情绪情感发育特点：}
\begin{enumerate}[leftmargin=*]
    \item 情绪的冲动性逐渐减少。
    \item 情绪情感以外显性为主，内隐性逐渐增强。
    \item 情绪情感以易变性为主，稳定性逐渐发展。
\end{enumerate}

\textbf{3. 学龄期儿童情绪情感发育特点：}
\begin{enumerate}[leftmargin=*]
    \item 情感内容逐渐丰富。
    \item 情感的深刻性和稳定性逐渐增加。
    \item 高级情感逐渐发展。
\end{enumerate}

\textbf{4. 青少年情绪情感发育特点：}
\begin{enumerate}[leftmargin=*]
    \item 情绪情感迅速而热烈，有两极性。
    \item 内隐性和表现性共同存在。
    \item 高级情感已经有了相当的发展。
\end{enumerate}
\end{keybox}

\section{儿童少年个性及社会化发展}

\subsection{个性的发展}

\begin{defbox}
\textbf{个性（personality）/ 人格：}个人的整个精神面貌，即具有一定倾向性的心理特征的总和，包括一个人个性倾向性、个性心理特征及自我意识。

\textbf{气质（temperament）：}婴幼儿期出生后最早表现出来的一种较为明显而稳定的个性特征，在婴幼儿社会性发展过程中有重要的地位和作用。
\end{defbox}

\subsection{自我意识}

\begin{defbox}
\textbf{自我意识（self-consciousness）：}由自我概念、自我评价和自我（情绪）体验等共同组成的心理过程，也是个体不断社会化的过程。

\begin{enumerate}[leftmargin=*]
    \item \textbf{自我概念（self-concept）：}个人心目中对自己的印象，包括对自己存在、个人身体能力、性格、态度、思想等方面的认识。
    \item \textbf{自我评价（self-evaluation）：}自我意识发展的主要成分和标志。
\end{enumerate}
\end{defbox}

\subsection{社会化}

\begin{defbox}
\textbf{社会化（socialization）：}个体在特定社会文化环境中，形成适合于该社会与文化的人格，掌握该社会公认的行为方式，成为合格社会成员的过程。
\end{defbox}

% ----- 复习题 -----
\reviewcontent{
\section{复习题}

\begin{enumerate}[leftmargin=*, itemsep=8pt]
    \item \textbf{【名词解释】}心理行为发育；认知；感知觉；感觉；知觉；注意；记忆；思维；想象；语言；前语言阶段；情绪；情感；个性/人格；气质；自我意识；自我概念；自我评价；社会化
    \item \textbf{【简答题】}简述儿童少年视觉、听觉的发育特点。
    \item \textbf{【简答题】}简述儿童少年思维发育的各阶段特点。
    \item \textbf{【简答题】}简述青少年情绪情感发育特点。
    \item \textbf{【简答题】}简述儿童少年个性及社会化发展的主要内容。
\end{enumerate}

\subsection*{参考答案要点}

\begin{enumerate}[leftmargin=*, itemsep=5pt]
    \item \textbf{认知}：大脑反映客观事物的特征、状态及其相互关系并揭示事物对人的意义与作用的高级心理活动。\textbf{感觉}：对事物个别属性的反映。\textbf{知觉}：人对事物整体属性的综合反映。\textbf{注意}：心理活动对一定对象有选择的集中，是认知过程的开始。\textbf{记忆}：人脑对过去经验的反映，包括识记、保持和再现。\textbf{思维}：人脑对客观事物间接、概括的反映，属认知的高级阶段。\textbf{情绪}：客观事物是否符合人的需要产生的态度体验。\textbf{情感}：与人的高级社会性需要满足与否相联系的态度体验。\textbf{个性}：个人的整个精神面貌，即具有一定倾向性的心理特征的总和。\textbf{自我意识}：由自我概念、自我评价和自我体验等共同组成的心理过程。\textbf{社会化}：个体在特定社会文化环境中形成适合于该社会与文化的人格，掌握该社会公认的行为方式，成为合格社会成员的过程。

    \item 视觉：(1)2个月视觉集中形成；(2)4个月出现辨色视觉；(3)6个月前视力发展敏感期；(4)2-3岁能辨别4种基本颜色并出现双眼视觉；(5)出生后3年内视觉系统基本成熟。听觉：(1)出生时听觉器官基本成熟；(2)2岁后听力接近成人；(3)学龄初期听觉感受性不断发展；(4)青少年期听觉能力基本稳定。

    \item (1)婴幼儿期：直觉行动思维，依靠感知觉和动作完成；(2)学龄前期：直观行动思维→具体形象思维→抽象逻辑思维，占主导的是具体形象思维；(3)学龄期：从具体形象思维逐步过渡到抽象逻辑思维（质的变化）；(4)青春期：抽象逻辑思维处于优势地位。

    \item (1)情绪情感迅速而热烈，有两极性；(2)内隐性和表现性共同存在；(3)高级情感已经有了相当的发展。

    \item 个性发展：个性是个人整个精神面貌，包括个性倾向性、个性心理特征及自我意识；气质是婴幼儿期最早表现出的明显而稳定的个性特征。自我意识：由自我概念（对自己存在、能力、性格等的认识）和自我评价（自我意识发展的主要成分和标志）组成。社会化：个体在特定社会文化环境中形成适合该社会与文化的人格，掌握公认行为方式的过程。
\end{enumerate}
}

\newpage
